\section{Background and Related Work}

\subsection{Retrieval-Augmented Generation}
Retrieval-Augmented Generation (RAG) is a common approach for building question-answering systems over external documents~\cite{lewis2020retrieval}. Instead of asking a language model to answer only from its internal knowledge, a RAG system first retrieves relevant passages from a document collection and then uses them as context for answer generation.

This approach is useful for legal question answering because legal information must come from official documents. However, a standard RAG pipeline is not always sufficient for law. If documents are divided into simple text chunks, the system may lose information about the document structure, such as articles, clauses, points, appendices, and legal rank. As a result, the retrieved passages may be textually relevant but legally incomplete.

In this thesis, RAG is used as the main framework for grounding answers in Vietnamese labor-law documents. The system retrieves legal evidence before generating an answer, and the final response is expected to remain connected to the retrieved legal sources.

\subsection{Legal Question Answering}
Legal question answering has stricter requirements than general question answering because legal answers must be traceable to legal authority and context~\cite{guha2023legalbench,zhong2020legalqa}. In many domains, an answer may be considered useful if it is fluent and semantically relevant. In law, this is not enough. A legal answer must also show the legal basis that supports it.

This requirement is especially important in Vietnamese labor law. A practical question may involve several related documents, such as the Labor Code, a decree, a circular, or an appendix. These documents may explain, limit, or guide each other. Some rules may also be replaced or clarified by later documents. Therefore, a legal QA system must not only find similar text, but also preserve the legal context around each provision.

For this reason, the system in this thesis treats citations as a core part of the answer. The answer should indicate the relevant legal document and provision, so that users can inspect where the information comes from.

\subsection{Hybrid Retrieval for Legal Evidence}
Retrieval is the first important step in a RAG system. Dense retrieval methods are useful because they search by semantic meaning rather than exact lexical overlap~\cite{karpukhin2020dense,reimers2019sentencebert}. They can find passages that are related to a user question even when the wording is different. However, legal questions often contain exact references, such as article numbers, clause numbers, document names, or specific legal terms. These details may be missed by semantic search alone.

Sparse or keyword-based retrieval is useful for exact matching of legal references, document names, and terminology~\cite{robertson2009probabilistic}. It can help find provisions that contain specific legal references or terms. For legal QA, both types of retrieval are useful: dense retrieval helps with meaning, while sparse retrieval helps with legal coordinates and exact wording.

This thesis therefore uses hybrid retrieval, combining semantic and lexical signals to support both meaning-based retrieval and exact legal-coordinate matching. The system combines semantic search with keyword-oriented search in Qdrant. This allows the retriever to handle both natural-language questions and questions that depend on specific legal references.

\subsection{Graph-Augmented Retrieval}
Legal documents are connected by structure and references. A law may contain general rules, while a decree provides implementation details. A circular may give forms, lists, or additional guidance. An appendix may contain the table or template needed to answer a specific question. These relationships are difficult to represent with a flat list of text chunks.

Graph-augmented retrieval addresses this problem by representing legal units and their relationships as a graph for retrieval support~\cite{edge2024graphrag}. In this thesis, Neo4j is used to store relationships between documents, articles, clauses, points, appendices, topics, and related legal provisions~\cite{neo4jdocs}. The graph helps the system expand from an initially retrieved provision to other provisions that may also be needed.

This is important because legal relevance is not always the same as textual similarity. A provision may be necessary for the answer because it is a parent rule, an implementing rule, a cross-reference, or an appendix, even if it does not use the same words as the user question.

\subsection{Evaluation Requirements for Legal QA}
Evaluating a legal QA system requires more than checking whether the answer sounds correct. A generated answer may be fluent but still cite the wrong article, miss an important provision, or answer a question that is outside the available corpus. Therefore, evaluation should check both retrieval and answer grounding, following the broader motivation of automated RAG evaluation frameworks~\cite{es2023ragas,saadfalcon2023ares}.

This is also consistent with legal NLP benchmark design, where legal understanding and reasoning are evaluated beyond general text fluency~\cite{chalkidis2021lexglue}.

In this thesis, retrieval is evaluated by whether the system finds the required legal citations and avoids distracting or forbidden citations. End-to-end answer quality is evaluated by checking whether the final answer is supported by retrieved evidence, whether its citations remain inside the retrieved context, and whether unsupported questions are refused or marked as insufficient context.

This evaluation design is intentionally deterministic. It focuses on reproducible software behavior rather than subjective judgment alone. Although this does not replace expert legal review, it provides a practical way to measure whether the system follows its main design goal: answering only when there is enough retrieved legal evidence.

\subsection{Research Gap}
Existing RAG systems show that external retrieval can improve the factual grounding of generated answers. Existing legal NLP research also shows that legal texts require careful treatment of context and citations. However, there remains a practical gap in building a deployed Vietnamese labor-law assistant that preserves legal hierarchy, retrieves related provisions, validates citations, and refuses unsupported questions within a scoped corpus.

This thesis addresses that gap by implementing a citation-grounded QA workflow for Vietnamese labor law. The focus is not on training a new language model, but on building a reliable software pipeline from legal corpus preparation to retrieval, graph-based evidence expansion, answer generation, citation validation, and benchmark evaluation.
